\section{Continuous Assessment 2 (CA-2) --- Set A}
\label{sec:ca2_set_a}

\LPUPaperHeader{Continuous Assessment (CA-2)}{A}{50 Minutes}{30 Marks}{CO4 (Multivariable Limits, Continuity, Partial Derivatives \& Optimization), CO5 (Multiple Integrals \& Transformations)}

\begin{center}
\sffamily\textbf{\color{AlertRed}Note: Attempt all six questions. Each question carries 5 marks.}
\end{center}

% ------------------------------------------------------------------------------
% QUESTION PAPER
% ------------------------------------------------------------------------------

\questionitem{1}{
    Find if the two-variable limit exists as $(x,y) \to (0,0)$:
    \[
    \lim_{(x,y)\to (0,0)} \frac{x^2 y}{x^4 + y^2}
    \]
    Show the rigorous mathematical justification using path testing.
}{5}{CO4}{L1 Remember / Analyze}

\questionitem{2}{
    For the function $f(x,y) = x^2 \tan^{-1}\left(\frac{y}{x}\right)$, verify the equality of mixed partial derivatives (Schwarz's Theorem):
    \[
    \frac{\partial^2 f}{\partial x \partial y} = \frac{\partial^2 f}{\partial y \partial x}
    \]
}{5}{CO4}{L3 Apply}

\questionitem{3}{
    A company manufactures two products $x$ and $y$. The joint profit function is given by:
    \[
    P(x,y) = 40x + 30y - (x^2 + xy + 2y^2)
    \]
    Find the production levels $x$ and $y$ that \textbf{maximize the profit}, and compute the maximum profit value.
}{5}{CO4}{L3 Apply}

\questionitem{4}{
    Change the order of integration and hence evaluate:
    \[
    I = \int_0^2 \int_{x-1}^1 (x - y) \, dy \, dx
    \]
}{5}{CO5}{L3 Apply}

\questionitem{5}{
    Evaluate the double integral by transforming to polar coordinates:
    \[
    \iint_R e^{x^2 + y^2} \, dA
    \]
    where $R$ is the circular disk of unit radius centered at the origin: $x^2 + y^2 \le 1$.
}{5}{CO5}{L3 Apply}

\questionitem{6}{
    Evaluate the triple integral:
    \[
    \int_0^1 \int_0^{\sqrt{1-x^2}} \int_0^1 x z \, dz \, dy \, dx
    \]
}{5}{CO5}{L3 Apply}

\vspace{5mm}
\hrule
\vspace{5mm}

% ------------------------------------------------------------------------------
% COMPLETE STEP-BY-STEP SOLUTIONS
% ------------------------------------------------------------------------------
\subsection*{Solutions \& Marking Scheme (CA-2 Set A)}

% Solution 1
\begin{solutionbox}[Solution to Question 1 (5 Marks)]
\textbf{Step 1: Path Testing along Linear Directions}
Consider paths approaching $(0,0)$ along straight lines $y = mx$:
\[
\lim_{x\to 0} \frac{x^2 (mx)}{x^4 + (mx)^2} = \lim_{x\to 0} \frac{m x^3}{x^2(x^2 + m^2)} = \lim_{x\to 0} \frac{mx}{x^2 + m^2} = \frac{0}{0 + m^2} = 0
\]
The limit along all straight lines is 0. However, this is NOT sufficient to prove that the limit exists!

\textbf{Step 2: Path Testing along Parabolic Curves}
Observe the powers in the denominator: $x^4 + y^2$. The terms have balanced degrees if $y^2 \propto x^4 \implies y = k x^2$.
Approach the origin along the parabola $y = k x^2$:
\[
\lim_{x\to 0} \frac{x^2 (k x^2)}{x^4 + (k x^2)^2} = \lim_{x\to 0} \frac{k x^4}{x^4 + k^2 x^4} = \lim_{x\to 0} \frac{k x^4}{x^4(1 + k^2)} = \mathbf{\frac{k}{1 + k^2}}
\]
\textbf{Step 3: Conclusion}
The limiting value depends explicitly on the parameter $k$ (the path of approach). For instance:
\begin{itemize}[leftmargin=6mm]
    \item Along $y = x^2$ ($k = 1$): limit is $\frac{1}{1 + 1} = \frac{1}{2}$.
    \item Along $y = 2x^2$ ($k = 2$): limit is $\frac{2}{1 + 4} = \frac{2}{5}$.
\end{itemize}
Since the limit along different approaching paths produces distinct values, the two-variable limit:
\[
\mathbf{\lim_{(x,y)\to (0,0)} \frac{x^2 y}{x^4 + y^2} \quad \text{DOES NOT EXIST}.}
\]

\begin{examinertip}[Exam Warning]
Never stop after checking $y = mx$. Always examine the degrees in the denominator. If the denominator is $x^{2a} + y^{2b}$, test paths of the form $y = k x^{a/b}$.
\end{examinertip}

\begin{markingrubric}
\textbf{Mark Distribution:}
$\bullet$ Testing linear paths $y = mx$ and finding 0: \textbf{1.5 Marks}\\
$\bullet$ Testing parabolic path $y = k x^2$ with balanced exponents: \textbf{2 Marks}\\
$\bullet$ Proving path dependence and concluding limit does not exist: \textbf{1.5 Marks}
\end{markingrubric}
\end{solutionbox}

% Solution 2
\begin{solutionbox}[Solution to Question 2 (5 Marks)]
Given $f(x,y) = x^2 \tan^{-1}\left(\frac{y}{x}\right)$.

\textbf{Step 1: Compute $\frac{\partial f}{\partial y}$}
Treat $x$ as constant:
\[
\frac{\partial f}{\partial y} = x^2 \cdot \frac{1}{1 + \left(\frac{y}{x}\right)^2} \cdot \frac{\partial}{\partial y}\left(\frac{y}{x}\right) = x^2 \cdot \frac{x^2}{x^2 + y^2} \cdot \frac{1}{x} = \frac{x^3}{x^2 + y^2}
\]
Now differentiate with respect to $x$ to find $\frac{\partial^2 f}{\partial x \partial y}$:
\[
\frac{\partial^2 f}{\partial x \partial y} = \frac{\partial}{\partial x}\left[\frac{x^3}{x^2 + y^2}\right] = \frac{(x^2 + y^2)\frac{d}{dx}(x^3) - x^3 \frac{d}{dx}(x^2 + y^2)}{(x^2 + y^2)^2}
\]
\[
= \frac{(x^2 + y^2)(3x^2) - x^3(2x)}{(x^2 + y^2)^2} = \frac{3x^4 + 3x^2 y^2 - 2x^4}{(x^2 + y^2)^2} = \mathbf{\frac{x^4 + 3x^2 y^2}{(x^2 + y^2)^2}}
\]

\textbf{Step 2: Compute $\frac{\partial f}{\partial x}$}
Using the product rule on $x^2 \cdot \tan^{-1}(y/x)$:
\[
\frac{\partial f}{\partial x} = 2x \tan^{-1}\left(\frac{y}{x}\right) + x^2 \cdot \frac{1}{1 + \left(\frac{y}{x}\right)^2} \left(-\frac{y}{x^2}\right) = 2x \tan^{-1}\left(\frac{y}{x}\right) - \frac{x^2 y}{x^2 + y^2}
\]
Now differentiate with respect to $y$ to find $\frac{\partial^2 f}{\partial y \partial x}$:
\[
\frac{\partial^2 f}{\partial y \partial x} = 2x \cdot \frac{1}{1 + (y/x)^2} \left(\frac{1}{x}\right) - \frac{\partial}{\partial y}\left[\frac{x^2 y}{x^2 + y^2}\right]
\]
First term:
\[
2x \cdot \frac{x^2}{x^2 + y^2} \cdot \frac{1}{x} = \frac{2x^2}{x^2 + y^2}
\]
Second term (quotient rule with respect to $y$, $x$ constant):
\[
\frac{\partial}{\partial y}\left[\frac{x^2 y}{x^2 + y^2}\right] = x^2 \left[\frac{(x^2 + y^2)(1) - y(2y)}{(x^2 + y^2)^2}\right] = \frac{x^2(x^2 - y^2)}{(x^2 + y^2)^2} = \frac{x^4 - x^2 y^2}{(x^2 + y^2)^2}
\]
Combine both terms over common denominator $(x^2+y^2)^2$:
\[
\frac{\partial^2 f}{\partial y \partial x} = \frac{2x^2(x^2 + y^2) - (x^4 - x^2 y^2)}{(x^2 + y^2)^2} = \frac{2x^4 + 2x^2 y^2 - x^4 + x^2 y^2}{(x^2 + y^2)^2} = \mathbf{\frac{x^4 + 3x^2 y^2}{(x^2 + y^2)^2}}
\]
\textbf{Step 3: Conclusion}
Both mixed partial derivatives are identical:
\[
\mathbf{\frac{\partial^2 f}{\partial x \partial y} = \frac{\partial^2 f}{\partial y \partial x} = \frac{x^4 + 3x^2 y^2}{(x^2 + y^2)^2}} \quad \text{(Verified!)}
\]

\begin{markingrubric}
\textbf{Mark Distribution:}
$\bullet$ First partial $\partial f/\partial y$ and mixed partial $\partial^2 f/\partial x \partial y$: \textbf{2 Marks}\\
$\bullet$ First partial $\partial f/\partial x$ and mixed partial $\partial^2 f/\partial y \partial x$: \textbf{2 Marks}\\
$\bullet$ Exact algebraic equivalence verification: \textbf{1 Mark}
\end{markingrubric}
\end{solutionbox}

% Solution 3
\begin{solutionbox}[Solution to Question 3 (5 Marks)]
\textbf{Step 1: First Partial Derivatives}
Profit function: $P(x,y) = 40x + 30y - x^2 - xy - 2y^2$.
\[
\frac{\partial P}{\partial x} = 40 - 2x - y
\]
\[
\frac{\partial P}{\partial y} = 30 - x - 4y
\]
\textbf{Step 2: Stationary (Critical) Points}
Set $\frac{\partial P}{\partial x} = 0$ and $\frac{\partial P}{\partial y} = 0$:
\begin{align}
2x + y &= 40 \implies y = 40 - 2x \label{eq:p1} \\
x + 4y &= 30 \label{eq:p2}
\end{align}
Substitute \eqref{eq:p1} into \eqref{eq:p2}:
\[
x + 4(40 - 2x) = 30 \implies x + 160 - 8x = 30 \implies -7x = -130 \implies x = \frac{130}{7} \approx 18.57
\]
Substitute $x$ back:
\[
y = 40 - 2\left(\frac{130}{7}\right) = \frac{280 - 260}{7} = \frac{20}{7} \approx 2.86
\]
Critical point is $\left(\frac{130}{7}, \frac{20}{7}\right)$.

\textbf{Step 3: Second Derivative Test (Hessian)}
Compute second-order partial derivatives:
\[
r = \frac{\partial^2 P}{\partial x^2} = -2
\]
\[
s = \frac{\partial^2 P}{\partial x \partial y} = -1
\]
\[
t = \frac{\partial^2 P}{\partial y^2} = -4
\]
Discriminant:
\[
rt - s^2 = (-2)(-4) - (-1)^2 = 8 - 1 = 7 > 0
\]
Since $rt - s^2 = 7 > 0$ and $r = -2 < 0$, the function achieves a **strict local maximum** at this critical point.

\textbf{Step 4: Maximum Profit Value}
\[
P_{\max} = 40\left(\frac{130}{7}\right) + 30\left(\frac{20}{7}\right) - \left[\left(\frac{130}{7}\right)^2 + \left(\frac{130}{7}\right)\left(\frac{20}{7}\right) + 2\left(\frac{20}{7}\right)^2\right]
\]
\[
= \frac{5200 + 600}{7} - \frac{16900 + 2600 + 800}{49} = \frac{5800}{7} - \frac{20300}{49} = \frac{40600 - 20300}{49} = \frac{20300}{49} = \mathbf{\frac{2900}{7} \approx 414.29}
\]
Optimal production levels: $\mathbf{x = \frac{130}{7}}$, $\mathbf{y = \frac{20}{7}}$.

\begin{markingrubric}
\textbf{Mark Distribution:}
$\bullet$ First partial derivatives and critical point equations: \textbf{1.5 Marks}\\
$\bullet$ Solving for $(x,y) = (130/7, 20/7)$: \textbf{1 Mark}\\
$\bullet$ Second derivative test ($rt - s^2 > 0$ and $r < 0$): \textbf{1.5 Marks}\\
$\bullet$ Maximum profit value calculation: \textbf{1 Mark}
\end{markingrubric}
\end{solutionbox}

% Solution 4
\begin{solutionbox}[Solution to Question 4 (5 Marks)]
\textbf{Step 1: Region of Integration}
Given integral:
\[
I = \int_0^2 \int_{x-1}^1 (x - y) \, dy \, dx
\]
The given limits correspond to:
\[
x \in [0, 2], \quad y \in [x-1, 1]
\]
The region is bounded by:
\begin{itemize}[leftmargin=6mm]
    \item Bottom curve: $y = x - 1 \implies x = y + 1$
    \item Top horizontal line: $y = 1$
    \item Vertical boundary: $x = 0$
\end{itemize}
Intersection points:
\begin{itemize}[leftmargin=6mm]
    \item $x = 0 \implies y = -1$ (at point $(0, -1)$)
    \item $x = 2 \implies y = 2 - 1 = 1$ (at point $(2, 1)$)
    \item At $y = 1$: $x = 0$ to $x = 2$.
\end{itemize}
The range of $y$ across the entire region is $y \in [-1, 1]$.

\textbf{Step 2: Reversing the Order (Horizontal Strips)}
For any horizontal strip at height $y$:
\begin{itemize}[leftmargin=6mm]
    \item $x$ enters at the left boundary: $x = 0$.
    \item $x$ exits at the right line: $x = y + 1$.
\end{itemize}
Thus, the limits in reversed order are:
\[
\mathbf{I = \int_{-1}^1 \int_0^{y+1} (x - y) \, dx \, dy}
\]

\textbf{Step 3: Evaluation}
Inner integral with respect to $x$:
\[
\int_0^{y+1} (x - y) \, dx = \left[ \frac{x^2}{2} - y x \right]_0^{y+1} = \frac{(y+1)^2}{2} - y(y+1) = \frac{y^2 + 2y + 1 - 2y^2 - 2y}{2} = \frac{1 - y^2}{2}
\]
Outer integral with respect to $y$:
\[
I = \int_{-1}^1 \frac{1 - y^2}{2} \, dy = 2 \int_0^1 \frac{1 - y^2}{2} \, dy = \int_0^1 (1 - y^2) \, dy = \left[ y - \frac{y^3}{3} \right]_0^1 = 1 - \frac{1}{3} = \mathbf{\frac{2}{3}}
\]

\begin{markingrubric}
\textbf{Mark Distribution:}
$\bullet$ Sketch/identification of original boundary lines and vertices: \textbf{1.5 Marks}\\
$\bullet$ Correctly setting new limits $y \in [-1, 1]$ and $x \in [0, y+1]$: \textbf{2 Marks}\\
$\bullet$ Accurate integration yielding final answer $2/3$: \textbf{1.5 Marks}
\end{markingrubric}
\end{solutionbox}

% Solution 5
\begin{solutionbox}[Solution to Question 5 (5 Marks)]
\textbf{Step 1: Polar Coordinate Transformation}
Let:
\[
x = r \cos\theta, \quad y = r \sin\theta \implies x^2 + y^2 = r^2
\]
The differential area element transforms as:
\[
dA = dx \, dy = r \, dr \, d\theta
\]
\textbf{Step 2: Region Limits}
The region $R$ is the complete unit disk $x^2 + y^2 \le 1$:
\[
r \in [0, 1], \quad \theta \in [0, 2\pi]
\]
\textbf{Step 3: Integral Evaluation}
\[
\iint_R e^{x^2+y^2} \, dA = \int_0^{2\pi} \int_0^1 e^{r^2} \cdot r \, dr \, d\theta
\]
Separate the independent variables:
\[
= \left( \int_0^{2\pi} d\theta \right) \left( \int_0^1 r e^{r^2} \, dr \right)
\]
The angular integral is $[ \theta ]_0^{2\pi} = 2\pi$.
For the radial integral, let $u = r^2 \implies du = 2r \, dr \implies r \, dr = \frac{du}{2}$:
\[
\int_0^1 r e^{r^2} \, dr = \frac{1}{2} \int_0^1 e^u \, du = \frac{1}{2} [e^u]_0^1 = \frac{1}{2}(e^1 - e^0) = \frac{e - 1}{2}
\]
Multiply the results:
\[
\iint_R e^{x^2+y^2} \, dA = 2\pi \left(\frac{e - 1}{2}\right) = \mathbf{\pi (e - 1)}
\]

\begin{markingrubric}
\textbf{Mark Distribution:}
$\bullet$ Transformation formulas and Jacobian $r \, dr \, d\theta$: \textbf{1.5 Marks}\\
$\bullet$ Correct limits for unit circle ($r \in [0,1], \theta \in [0, 2\pi]$): \textbf{1.5 Marks}\\
$\bullet$ Integration and final exact answer $\pi(e - 1)$: \textbf{2 Marks}
\end{markingrubric}
\end{solutionbox}

% Solution 6
\begin{solutionbox}[Solution to Question 6 (5 Marks)]
\textbf{Step 1: Understand the Integral Structure}
\[
I = \int_0^1 \int_0^{\sqrt{1-x^2}} \int_0^1 x z \, dz \, dy \, dx
\]
Notice that the innermost integral depends on $z$, with constant limits $z \in [0, 1]$.
\textbf{Step 2: Inner Integration (with respect to $z$)}
\[
\int_0^1 x z \, dz = x \left[ \frac{z^2}{2} \right]_0^1 = \frac{1}{2} x
\]
\textbf{Step 3: Middle Integration (with respect to $y$)}
Now the integral becomes:
\[
\int_0^{\sqrt{1-x^2}} \frac{1}{2} x \, dy = \frac{1}{2} x [y]_0^{\sqrt{1-x^2}} = \frac{1}{2} x \sqrt{1 - x^2}
\]
\textbf{Step 4: Outermost Integration (with respect to $x$)}
\[
I = \int_0^1 \frac{1}{2} x \sqrt{1 - x^2} \, dx
\]
Let $u = 1 - x^2 \implies du = -2x \, dx \implies x \, dx = -\frac{du}{2}$.
When $x = 0, u = 1$; when $x = 1, u = 0$:
\[
I = \frac{1}{2} \int_1^0 \sqrt{u} \left(-\frac{du}{2}\right) = \frac{1}{4} \int_0^1 u^{1/2} \, du = \frac{1}{4} \left[ \frac{u^{3/2}}{3/2} \right]_0^1 = \frac{1}{4} \cdot \frac{2}{3} = \mathbf{\frac{1}{6}}
\]

\begin{markingrubric}
\textbf{Mark Distribution:}
$\bullet$ Integration with respect to $z$: \textbf{1.5 Marks}\\
$\bullet$ Integration with respect to $y$: \textbf{1.5 Marks}\\
$\bullet$ Substitution and final evaluation with respect to $x$ giving $1/6$: \textbf{2 Marks}
\end{markingrubric}
\end{solutionbox}

\newpage
